% ЗӨВХӨН БҮТЭН НООРОГТ (main.tex) хэвлэгдэнэ. Багшид үзүүлэх ноорогт (main_teacher.tex) орохгүй.
% 2.4 Төсөөтэй системүүд, 2.5 Шинэлэг тал — \iffalse-ээр түр унтраалттай (11/11 хэрэглэгч)
% 2026.11.11 (2), хэрэглэгч: 2.4 «Төсөөтэй системүүдийн харьцуулалт» ба 2.5 «Шинэлэг тал»-ыг
% түр унтраасан (ажил дөнгөж эхэлж байгаа тул одоохондоо их). Устгаагүй: буцааж асаахдаа
% доорх \iffalse ба төгсгөлийн \fi хоёр мөрийг арилгана. «Хаанаас авах» нь хүн ба байршил
% гэдгийг, мэдээллийн эх нь сайт ба сургуулийн албад гэдгийг тэр үед хүснэгтэд тулгана.
\iffalse
\section{Төсөөтэй системүүдийн харьцуулалт}
\label{sec:lit-compare}

Сүүлийн жилүүдэд их сургуулиуд оюутны асуултад хариулдаг туслах систем хийж байна. \tref{tab:similar_systems}-т дөрвөн системийг манай төлөвлөж буй системтэй харьцуулав.

\begin{table}[H]
\centering
\caption{Төсөөтэй системүүдийн харьцуулалт}
\label{tab:similar_systems}
\small
\begin{tabular}{|p{2.2cm}|p{1.9cm}|p{1.6cm}|p{3.0cm}|p{2.4cm}|p{2.4cm}|}
\hline
\textbf{Систем} & \textbf{Хэл} & \textbf{Асуух} & \textbf{Мэдээллийн эх} & \textbf{Хариулт} & \textbf{Үнэлгээ} \\
\hline
Unimib \cite{antico2024unimib} & Итали & Бичгээр & Баримт, холбоос & ChatGPT бичнэ & 6 оюутны туршилт \\
\hline
Marcel \cite{trienes2025marcel} & Герман & Бичгээр & Сайтын 249 баримт, 36 бэлэн хариулт & Загвар бичнэ, эхгүй бол татгалзана & 95 асуулт \\
\hline
ALOHA \cite{tao2025aloha} & Хятад, англи, франц & Бичгээр & Гарын авлага, сайт, хүснэгт & GPT-4o бичнэ & Бодит 100 асуулт \\
\hline
Burch \cite{matic2026burch} & Англи & Бичгээр & Сайтын 470+ хуудас & Загвар бичнэ & 20 асуулт \\
\hline
\textbf{Манай (төлөвлөж буй)} & \textbf{Монгол} & \textbf{Дуугаар, бичгээр} & \textbf{Сайт ба сургалтын албаны хүснэгт} & \textbf{Эхийн үгээр} & \textbf{Дуу таних: 12 бичлэг; хайлт: бодит асуулт цуглуулж байна} \\
\hline
\end{tabular}
\end{table}

Хүснэгтийн системүүд бүгд бичгээр асуудаг бөгөөд ихэнх нь хариултыг хэлний загвараар шинээр бичүүлдэг. Unimib-ийн оюутнууд мэдээлэл үргэлж зөв биш, заримдаа байхгүй холбоос зохиосныг дурдсан \cite{antico2024unimib}. Манай системийг эхийн үгийг өөрчлөхгүй өгдөг байхаар төлөвлөж байна.

\section{Судалгааны ажлын шинэлэг тал}
\label{sec:lit-novelty}

Бид 2026 оны 11-р сард вэб хайлтаар өмнөх ажлыг хайсан. Энэ нь системтэй тойм биш. Иймээс доорх зүйлсийг «огт байхгүй» гэж биш, бидний хайлтад олдоогүй хослол гэж үзнэ.
\begin{enumerate}
    \item Монгол хэлээр, сургуулийн ресепшн ба гар утсанд зориулсан, дуугаар болон бичгээр асуух боломжтой туслах.
    \item Онлайн яриа таних үйлчилгээ ба Монгол хэлээр сургасан Whisper-ийг сургуулийн асуултын ижил бичлэг дээр харьцуулсан. Энэ нь нэг хүний бичлэгээр хийсэн жижиг хэмжилт.
    \item Гараар хийсэн мэдлэгийн модны хүрээгээр хайлтад нэмэлт оноо өгөх аргыг Монгол сургуулийн бодит асуултаар шалгах.
\end{enumerate}
\fi
